\section{Upper-edge two-saddle approximation with growing cap}\label{sec:growing-cap-upper-saddle}

2026-11-26 \enspace Private research note

There are absolute constants $h_0>0$ and $C_*>0$ with the following property. Let $d\to\infty$, $1\le r<d$, and $n\ge C_*d^4\log(d+1)$. Take $t>0$ and let $S$ have probabilities proportional to $\binom ds t^s$ for $0\le s\le r$. With $J=r-S$, assume

\begin{equation*}
h=\mathbb E J\le h_0,\qquad nh/d^2\longrightarrow\infty.
\end{equation*}

Let $H(w)=\sum_{s\le r}\binom ds w^s$, $\phi=\pi/d$, $H_\tau=H(te^{i\tau})$, and

\begin{equation*}
\lambda_s=\frac{s(d-s)}{d(d-1)},\qquad
\Lambda_\tau=\frac{\mathbb E(\lambda_S e^{-i\tau J})}{\mathbb E e^{-i\tau J}},
\qquad \Lambda_0=\mathbb E\lambda_S.
\end{equation*}

For the balanced-array polynomial $P_n$, the uniform approximation in this regime is

\begin{equation}
P_n(-t^d)=2\mathrm{Re}\left(
 \frac{H_\phi^n \Lambda_\phi^{-(d-1)/2}}{\sqrt d(2\pi n)^{(d-1)/2}}
 \right)
 +o\left(\frac{|H_\phi|^n|\Lambda_\phi|^{-(d-1)/2}}{\sqrt d(2\pi n)^{(d-1)/2}}\right).
\tag{S1}
\end{equation}

The power uses the logarithm in the right half-plane. The error is absolute relative to the magnitude of one saddle, including when the main real part is small. The condition $nh/d^2\to\infty$ may hold arbitrarily slowly.

\subsection{Tilt and diagonal gaps}

The ratio of successive unnormalized deficit weights at $j,j+1$ is $(r-j)/[(d-r+j+1)t]$, decreasing in $j$. The first ratio $q$ satisfies $q=\Pr(J=1)/\Pr(J=0)\le h/(1-h)$. Thus $\Pr(J=j)\le q^j$ and $\mathbb E J^\ell=O_\ell(h)$ for every fixed positive integer $\ell$, uniformly once $h_0$ is sufficiently small. Also $h\le Cq$ by the same bound, so $\Pr(J=1)\asymp h$ and $\mathrm{Var}(J)\asymp h$. For $j\le r$,

\begin{equation*}
\frac{\lambda_{r-j}}{\lambda_r}
=\frac{r-j}{r}\frac{d-r+j}{d-r}\le1+j.
\end{equation*}

As $\Pr(J=0)\ge1-h$, this implies $\Lambda_0\asymp\lambda_r$, $c/d\le\Lambda_0\le C$, and $\mathbb E[\lambda_S^a J^\ell]=O_{a,\ell}(\Lambda_0^a h)$ for fixed $a,\ell\ge1$. The same bound holds for the expectation of $\lambda_S^a\mathbf1_{J>0}$.

Write $g(\tau)=\mathbb E e^{-i\tau J}$, so $H_\tau/H(t)=e^{ir\tau}g(\tau)$. Uniformly on $[-\pi,\pi]$, $g(\tau)=1+O(h)$ and $|g(\tau)|\le\exp[-ch(1-\cos\tau)]$, the latter by retaining the $J=0,1$ terms in the independent-copy formula for $1-|g|^2$. Centering $J$ and differentiating its characteristic function gives $(\log|g(\tau)|)''=-\mathrm{Var}(J)+O(h|\tau|)$ near zero by the centered third-moment bound. Evenness gives $\log|g(\tau)/g(\phi)|\le-ch(\tau^2-\phi^2)$ for $\phi\le|\tau|\le\eta_0$ with fixed sufficiently small $\eta_0$, and $|g(\phi)|\ge e^{-Ch/d^2}$. Beyond $\eta_0$, the global diagonal gap and the lower bound at $\phi$ give the same comparison with another constant. Thus

\begin{equation}
\log|H_\tau/H_\phi|\le-ch(\tau^2-\phi^2)
\qquad(\phi\le|\tau|\le\pi).
\tag{S2}
\end{equation}

For later volume comparisons, differentiating the numerator and denominator of $\Lambda_\tau$ once or twice costs $O(\Lambda_0h)$ and $O(h)$, respectively. This uses the $J$-moment bounds and $|g(\tau)|\ge1-O(h)$. Uniformly, $\Lambda_\tau=\Lambda_0(1+O(h))$, as follows by separating the $J=0$ term in numerator and denominator. Conjugation makes the first derivative at zero purely imaginary. Hence

\begin{equation}
\mathrm{Re}\Lambda_\tau/\Lambda_0=1+O(h\min(\tau^2,1)),\qquad
|\mathrm{Im}\Lambda_\tau|/\Lambda_0=O(h\min(|\tau|,1)).
\tag{S3}
\end{equation}

In particular $|\Lambda_\phi|/\Lambda_0=1+O(h/d^2)$.

\subsection{Localization relative to the principal modulus}

Let $R(\mathbf z)=\sum_{|A|\le r}t^{|A|}\prod_{i\in A}z_i$ on the torus with $\prod_i z_i=-1$, with normalized Haar measure. The integral of $R^n$ is $P_n(-t^d)$ by character orthogonality. Put $\overline z=d^{-1}\sum_i z_i$ and $D_z=\sum_i|z_i-\overline z|^2$.

For fixed size $s$ let $m=\min(s,d-s)$. Generate a uniform subset by putting $m$ disjoint random pairs in the first $2m$ positions of a uniform permutation, choosing one element from each pair by independent fair coins, and including the next $s-m$ positions. By symmetry the subset is uniform. Conditional on the permutation the modulus of its Fourier response is at most $e^{-X/8}$ where $X$ is the sum of the squared chords within the pairs. Indeed the squared modulus of the mean of the two points of a pair is one minus one fourth of the squared chord. For one ordered random pair the mean squared chord is $2D_z/(d-1)$, so $\mathbb E X=2mD_z/(d-1)$. Each squared chord is at most four. For two disjoint pairs their product expectation is at most $d(d-1)/[(d-2)(d-3)]$ times the squared one-pair expectation, by discarding the disjointness restriction in the numerator before dividing by the number of disjoint pairs. Thus $\mathbb E X^2\le C(\mathbb EX+(\mathbb EX)^2)$. From Cauchy--Schwarz,

\begin{equation*}
\mathbb E\frac{X}{1+X}\ge
\frac{(\mathbb E X)^2}{\mathbb E(X+X^2)}.
\end{equation*}

Since $1-e^{-x/8}\ge c x/(1+x)$, the fixed-size response has modulus at most $1-c\min(mD_z/d,1)$. This also holds for $m=0$. Averaging over $S$ and using $m\ge s(d-s)/d$ gives

\begin{equation}
|R/H(t)|\le1-c\Lambda_0\min(D_z,1).
\tag{S4}
\end{equation}

Here $\min(mD_z/d,1)\ge(m/d)\min(D_z,1)$ and $\mathbb E[S(d-S)/d]=(d-1)\Lambda_0$.

For $D_z\ge1/d$, (S4) gives modulus at most $H(t)(1-c/d^2)$. By reducing $h_0$ this is at most $|H_\phi|(1-c'/d^2)$, since $|H_\phi|/H(t)\ge e^{-Ch/d^2}$. Its absolute integral, divided by the principal scale

\begin{equation*}
|H_\phi|^n I_0,\qquad I_0=(2\pi n\Lambda_0)^{-(d-1)/2}/\sqrt d,
\end{equation*}

is bounded by $\exp(-c'n/d^2+O(d\log n))$, which is negligible in the asserted range.

For $D_z<1/d$, choose a pin minimizing $\sum_i|z_i-z_j|^2$, which is at most $2D_z$ by averaging. Lift the angles relative to this pin, take their average $\tau$ modulo $2\pi$ in $(-\pi,\pi]$, and write $z_i=e^{i(\tau+y_i)}$ with $\sum_i y_i=0$ and $\|y\|^2=O(1/d)$. The norm bound follows from the small-chord bound and subtracting the mean of the lifted differences. The common phase is an odd multiple of $\phi$ modulo $2\pi$. The representation is unique within this small norm domain: if two lifts give the same point, their relative shifts by integer multiples of $2\pi$ agree at all coordinates because the pairwise angular differences are small, and averaging fixes the canonical phase. In free coordinates the measure is $(2\pi)^{-(d-1)}dy_1\cdots dy_{d-1}$, with $y_d=-\sum_{i<d}y_i$. Each such chart contains the full ball $\|y\|\le d^{-1}$, because $D_z\le\sum_i|z_i-e^{i\tau}|^2\le \|y\|^2$.

\subsection{Transverse expansion at all common phases}

For $L=\sum_{i\in A}y_i$ with $A$ a uniform size-$s$ subset, write $P_j=\sum_i y_i^j$ and $m_j=(s)_j/(d)_j$. Conditional means vanish. Expansion using $P_1=0$ gives

\begin{equation*}
\begin{aligned}
\mathbb E(L^2\mid S=s)&=\lambda_s P_2,\\
\mathbb E(L^3\mid S=s)&=\lambda_s\frac{d-2s}{d-2}P_3,\\
\mathbb E(L^4\mid S=s)&=(m_1-7m_2+12m_3-6m_4)P_4
                  +3(m_2-2m_3+m_4)P_2^2.
\end{aligned}
\end{equation*}

For the cubic, the mixed sum $\sum_{i\ne j}y_i^2y_j$ is $-P_3$ and $e_3(\mathbf y)=P_3/3$. For the quartic the sums with respective patterns $3+1$, $2+2$, $2+1+1$, and $1+1+1+1$, with equal exponents unlabeled, are $-P_4$, $(P_2^2-P_4)/2$, $P_4-P_2^2/2$, and $P_2^2/8-P_4/4$. The coefficient of $P_4$ is $O(\lambda_s)$ by writing it in terms of $m_j-m_{j+1}=m_j(d-s)/(d-j)=O(\lambda_s)$. The coefficient of $P_2^2$ is $O(\lambda_s^2)$ by $m_2-2m_3+m_4=(s)_2(d-s)_2/(d)_4$. In particular $\mathbb E(L^4\mid S=s)\le C\lambda_s\|y\|^4$. Cauchy--Schwarz against the conditional second moment gives $\mathbb E(|L|^3\mid S=s)\le C\lambda_s\|y\|^3$.

Normalize the complex size weights by $g(\tau)$ and tilt by $e^{-i\tau J}$. Its absolute weights are bounded by a constant times the underlying real probabilities, uniformly on the entire circle. The conditional linear mean vanishes. Taylor's bound for the real exponential through degree two gives

\begin{equation}
\log(R/H_\tau)=-\tfrac12\Lambda_\tau\|y\|^2+O(\Lambda_0\|y\|^3)
\tag{S5}
\end{equation}

throughout the localized region. In detail, the exponential remainder is bounded by the conditional third absolute moment, and $|R/H_\tau-1|=O(\Lambda_0\|y\|^2)$ there by the vanishing linear term. The error from taking the logarithm is thus absorbed in (S5), with the branch near one. By (S3), $\mathrm{Re}\Lambda_\tau\ge(1-Ch_0)\Lambda_0$ for all phases. For small fixed $h_0$, (S5) gives real transverse decay at least $c\Lambda_0\|y\|^2$. Thus, in the retained region but outside the balls of radius $d^{-1}$ about all common phases, $|R|\le|H_\phi| e^{-c/d^3}$ by (S2). The absolute integral there relative to $|H_\phi|^n I_0$ is at most $\exp[-cn/d^3+O(d\log n)]$. At $n=C_* d^4\log(d+1)$, the ratio $n/(d^4\log n)$ tends to $C_*/4$. This ratio increases with $n$ in the stated region. Taking $C_*$ sufficiently large makes the last error negligible uniformly. The same choice also covers the earlier global tail.

\subsection{Summing the phase balls and evaluating the principal integrals}

Inside a ball $\|y\|\le d^{-1}$, (S3) and (S5) bound the absolute integrand by a Gaussian with quadratic coefficient $\Lambda_0[1-O(d^{-1}+h\min(\tau^2,1))]$. Its integral relative to $I_0$ is at most $\exp[O(1+dh\min(\tau^2,1))]$. The zero-sum Gaussian determinant in the free coordinates is $d$, since the quadratic form matrix is $I+\mathbf1\mathbf1^\top$. For every nonprincipal phase $|\tau|\ge3\phi$ the factor involving $dh\tau^2$ is absorbed by half the gap in (S2), because $\tau^2\le(9/8)(\tau^2-\phi^2)$ and $n\gg d$. Summing the odd phase labels gives relative error at most

\begin{equation*}
C\sum_{j\ge1}\exp[-c(nh/d^2)((2j+1)^2-1)]
=O(e^{-c'nh/d^2})=o(1).
\end{equation*}

There is no prefactor growing with $d$. This is the reason the condition $nh/d^2\to\infty$ can be arbitrarily slow.

At $\tau=\phi$, compare $\exp(n\log(R/H_\phi))$ with $\exp(-n\Lambda_\phi\|y\|^2/2)$ by $|e^A-e^B|\le |A-B|e^{\max(\mathrm{Re} A,\mathrm{Re} B)}$. The weakened Gaussian just used has normalization $O(I_0)$ at $\phi$. Its normalized third norm moment is $O(d^{3/2}/(n\Lambda_0)^{3/2})$. Consequently the integrated absolute error of (S5) is at most

\begin{equation*}
O(I_0\,d^{3/2}/\sqrt{n\Lambda_0})
=O(I_0\,d^2/\sqrt n)
=o(I_0).
\end{equation*}

The comparison Gaussian extends to the full hyperplane at cost $o(I_0)$, since its absolute tail outside $d^{-1}$ is bounded relative to $I_0$ by $C2^{d/2}\exp(-c n\Lambda_0/d^2)$ using the Gaussian quadratic exponential moment. Its complex integral is $(2\pi n\Lambda_\phi)^{-(d-1)/2}/\sqrt d$ and its modulus is $(1+o(1))I_0$ by (S3). The ball at $-\phi$ gives the conjugate integral. This completes (S1).
